\section{Introduction}
\label{sec:intro}

Fire in the UK residential built environment remains a large and, recently, a
growing problem: fire and rescue services attended 40{,}350 building fires in
England in the year ending December 2025, up 5.5\% year-on-year, of which 26{,}298
were dwelling fires \citep{mhclg2026firestats}. Yet the fire risk assessment (FRA)
that governs those buildings is still periodic, qualitative, and compliance-driven
--- checklist judgements recorded under the Regulatory Reform (Fire Safety)
Order \citep{rrfso2005} and codified methodologies such as PAS~79
\citep{pas79_1_2020,pas79_2_2020}. The Grenfell Tower fire and the Phase~2 Inquiry
reframed residential fire safety around evidence, accountability, and the active
management of higher-risk buildings \citep{grenfell2024phase2}, and the regulatory
response --- the Fire Safety Act 2021 \citep{firesafetyact2021} and the Building
Safety Act 2022 \citep{bsa2022}, with their safety-case regime --- demands exactly
the kind of quantitative, evidence-updated risk picture that current practice does
not produce.

Three gaps stand out. Risk is communicated as a categorical score rather than a
distribution, despite evidence that risk matrices have poor resolution and can rank
worse than random \citep{cox2008,thomas2014}. Inspection findings are treated as
ground truth rather than as noisy, intermittent observations; a targeted search
across the fire-risk-assessment, inter-rater-reliability, assessor-variability, and
PAS~79 literatures surfaced competence standards and professional guidance, but we
are not aware of any empirical study estimating inter-assessor variability from
repeated assessments of the same UK residential building --- the sector's response
has been a competence standard \citep{bs8674_2025} rather than measurement
of assessor noise --- motivating the noisy-observation treatment we import from
structural risk-based inspection \citep{straub2005,luque2019}. And fire outcomes are
rare events, yet calibration and uncertainty quantification (UQ) are seldom made
explicit in operational FRA, even as reviews of fire machine learning confirm that
UQ and calibration remain underdeveloped \citep{tapehnaser2022} against an established
calibration-and-sharpness standard \citep{gneiting2007}.

We recast FRA as Bayesian inference under uncertainty: a building's fire risk is a
posterior distribution over ignition, spread, evacuation failure, and consequence
severity, inferred from heterogeneous public evidence and updateable as new
observations arrive. The novelty is not a new Bayesian network, a new fire
simulator, or a new evacuation model in isolation; it is an auditable decision
framework that couples calibrated UK open-data rare-event modelling,
mechanistically grounded consequence simulation, uncertainty propagation, and
decision-theoretic prevention targeting into one posterior-risk workflow. The paper
is explicit throughout about what is built and what is specified: the calibrated
public-data and simulation components below are \emph{implemented}, while the
latent-state, noisy-inspection extension for live risk updating
(Section~\ref{sec:model}) is \emph{proposed} --- fully specified but not yet fully
instantiated. Five contributions follow:

\begin{enumerate}
  \item A UK open-data fire-risk database linking incident, building-stock,
        deprivation, occupancy, and vulnerability evidence at dwelling and
        small-area scale, validated against published totals
        (Section~\ref{sec:data}).
  \item A calibrated rare-event occurrence model for dwelling fires ---
        hierarchical negative-binomial with out-of-time probabilistic validation
        (Section~\ref{sec:results}).
  \item Consequence models that explicitly separate predictive association from
        intervention effect, with formal treatment of alarm-state confounding,
        alongside a mechanism-informed archetype fire/evacuation simulator that
        supplies consequence estimates where public geometry and evacuation data
        do not exist (Section~\ref{sec:surrogate}).
  \item A decision-theoretic targeting demonstration: under stated assumptions, a
        fixed prevention budget allocated by posterior risk and expected benefit
        returns roughly nine times the value of untargeted allocation and 2.5
        times deprivation-only targeting (Section~\ref{sec:results-decision}).
  \item A proposed latent-state formulation in which inspection findings enter as
        noisy observations (Section~\ref{sec:model}), motivated by the
        measurement-error tradition of structural risk-based inspection, together
        with a real-building stress test bounding the simulator's regime
        assumptions against Grenfell Tower
        (Section~\ref{sec:results-casestudy}).
\end{enumerate}
